\section{Approximate cell decomposition}
\label{sec:l13-approximate}
Exact decomposition follows obstacle geometry. An alternative is to cover
configuration space with simple cells and refine the cells whose status is
uncertain. In two dimensions, a square can be divided into four equal
squares; repeating this operation produces a \emph{quadtree}. Small cells
resolve obstacle boundaries while large cells can represent open regions.

Consider a bounded two-dimensional configuration space. Enclose it in a
square if necessary, treating the region outside the permitted domain as
forbidden. Classify each cell by its interior:
\begin{itemize}
\item \textbf{Empty:} its interior contains only free configurations.
\item \textbf{Occupied:} its interior is entirely forbidden.
\item \textbf{Mixed:} it contains both free and forbidden configurations.
\end{itemize}
Obstacle contact still counts as collision. Calling a cell empty does
not make every point on its boundary free: shared boundaries must be
checked when constructing transitions.

\subsection{Refinement guided by a connectivity graph}
The following procedure uses an \emph{optimistic} graph: mixed cells are
temporarily allowed, but a route through them is only a candidate.
\begin{enumerate}
\item Split the initial square into four cells and classify them.
\item Create a graph node for each nonoccupied cell. Connect cells that
share a boundary segment of positive length through which a transition
might be possible. Do not connect cells merely because they share a corner.
For two empty cells, require a free portion of their common boundary.
\item Locate the cells containing $q_S$ and $q_G$ and search the graph
between them. First check that the start and goal themselves are free.
Prefer a route using only empty cells: search that subgraph first, and
use the optimistic graph if no such route is available yet.
\item If the candidate route includes mixed cells, subdivide mixed cells
into four, update the graph, and search again. A simple implementation
refines all mixed cells at each iteration.
\item Once the route uses only empty cells, choose a free crossing point
on each shared boundary and construct a continuous path through those
points. Validate the resulting path before accepting it.
\end{enumerate}
A conservative alternative excludes mixed cells from the graph. In that
version, an unsuccessful search can trigger further subdivision: a passage
may exist inside cells that have not yet been resolved.

Figure~\ref{fig:l13-quadtree} illustrates three successive resolutions.
The initial graph is small but cannot certify a route. Refinement reveals
occupied cells inside the larger obstacle and creates smaller cells around
both obstacle boundaries. A large empty square remains in the lower-left
region throughout: there is no need to subdivide every cell uniformly.

\begin{figure}[p]
\centering% Generated by tools/generate_sampling_figures.py
\begin{tikzpicture}[scale=0.65,line cap=round,line join=round,
 every node/.style={font=\footnotesize},>=Stealth]
\begin{scope}[yshift=-0.0cm]
\node[anchor=west,font=\small\bfseries] at (0,8.6) {Iteration 1: 4 cells};
\fill[white] (0.0000,0.0000) rectangle (4.0000,4.0000);
\fill[black!12] (4.0000,0.0000) rectangle (8.0000,4.0000);
\fill[black!12] (0.0000,4.0000) rectangle (4.0000,8.0000);
\fill[black!12] (4.0000,4.0000) rectangle (8.0000,8.0000);
\fill[black!85] (1.8000,5.2000) -- (2.7000,7.0000) -- (3.7000,5.2000) -- cycle;
\fill[black!85] (4.0000,2.0000) -- (6.0000,2.0000) -- (6.0000,2.5000) -- (6.8000,2.5000) -- (6.8000,6.7000) -- (5.0000,6.7000) -- (5.0000,6.0000) -- (4.0000,6.0000) -- cycle;
\draw[black!45,thin] (0.0000,0.0000) rectangle (4.0000,4.0000);
\draw[black!45,thin] (4.0000,0.0000) rectangle (8.0000,4.0000);
\draw[black!45,thin] (0.0000,4.0000) rectangle (4.0000,8.0000);
\draw[black!45,thin] (4.0000,4.0000) rectangle (8.0000,8.0000);
\draw[thick] (0,0) rectangle (8,8);
\fill[robotframe] (0.4000,7.3000) circle (2pt);
\node[fill=white,inner sep=1pt,above right] at (0.4000,7.3000) {$q_S$};
\fill[robotframe] (7.5000,1.3000) circle (2pt);
\node[fill=white,inner sep=1pt,above right] at (7.5000,1.3000) {$q_G$};
\begin{scope}[xshift=9.5cm]
\draw[black!20] (0,0) rectangle (8,8);
\draw[black!50] (6.0000,2.0000) -- (2.0000,2.0000);
\draw[black!50] (2.0000,6.0000) -- (2.0000,2.0000);
\draw[black!50] (6.0000,6.0000) -- (6.0000,2.0000);
\draw[black!50] (6.0000,6.0000) -- (2.0000,6.0000);
\draw[robotframe,very thick] (2.0000,6.0000) -- (2.0000,2.0000) -- (6.0000,2.0000);
\draw[fill=white] (2.0000,2.0000) circle (3pt);
\draw[fill=black!40] (6.0000,2.0000) circle (3pt);
\draw[fill=black!40] (2.0000,6.0000) circle (3pt);
\draw[fill=black!40] (6.0000,6.0000) circle (3pt);
\node[above left,fill=white,inner sep=1pt] at (2.0000,6.0000) {$C_S$};
\node[above left,fill=white,inner sep=1pt] at (6.0000,2.0000) {$C_G$};
\end{scope}
\end{scope}
\begin{scope}[yshift=-9.4cm]
\node[anchor=west,font=\small\bfseries] at (0,8.6) {Iteration 2: 13 cells};
\fill[white] (0.0000,0.0000) rectangle (4.0000,4.0000);
\fill[white] (4.0000,0.0000) rectangle (6.0000,2.0000);
\fill[white] (6.0000,0.0000) rectangle (8.0000,2.0000);
\fill[black!65] (4.0000,2.0000) rectangle (6.0000,4.0000);
\fill[black!12] (6.0000,2.0000) rectangle (8.0000,4.0000);
\fill[black!12] (0.0000,4.0000) rectangle (2.0000,6.0000);
\fill[black!12] (2.0000,4.0000) rectangle (4.0000,6.0000);
\fill[white] (0.0000,6.0000) rectangle (2.0000,8.0000);
\fill[black!12] (2.0000,6.0000) rectangle (4.0000,8.0000);
\fill[black!65] (4.0000,4.0000) rectangle (6.0000,6.0000);
\fill[black!12] (6.0000,4.0000) rectangle (8.0000,6.0000);
\fill[black!12] (4.0000,6.0000) rectangle (6.0000,8.0000);
\fill[black!12] (6.0000,6.0000) rectangle (8.0000,8.0000);
\fill[black!85] (1.8000,5.2000) -- (2.7000,7.0000) -- (3.7000,5.2000) -- cycle;
\fill[black!85] (4.0000,2.0000) -- (6.0000,2.0000) -- (6.0000,2.5000) -- (6.8000,2.5000) -- (6.8000,6.7000) -- (5.0000,6.7000) -- (5.0000,6.0000) -- (4.0000,6.0000) -- cycle;
\draw[black!45,thin] (0.0000,0.0000) rectangle (4.0000,4.0000);
\draw[black!45,thin] (4.0000,0.0000) rectangle (6.0000,2.0000);
\draw[black!45,thin] (6.0000,0.0000) rectangle (8.0000,2.0000);
\draw[black!45,thin] (4.0000,2.0000) rectangle (6.0000,4.0000);
\draw[black!45,thin] (6.0000,2.0000) rectangle (8.0000,4.0000);
\draw[black!45,thin] (0.0000,4.0000) rectangle (2.0000,6.0000);
\draw[black!45,thin] (2.0000,4.0000) rectangle (4.0000,6.0000);
\draw[black!45,thin] (0.0000,6.0000) rectangle (2.0000,8.0000);
\draw[black!45,thin] (2.0000,6.0000) rectangle (4.0000,8.0000);
\draw[black!45,thin] (4.0000,4.0000) rectangle (6.0000,6.0000);
\draw[black!45,thin] (6.0000,4.0000) rectangle (8.0000,6.0000);
\draw[black!45,thin] (4.0000,6.0000) rectangle (6.0000,8.0000);
\draw[black!45,thin] (6.0000,6.0000) rectangle (8.0000,8.0000);
\draw[thick] (0,0) rectangle (8,8);
\fill[robotframe] (0.4000,7.3000) circle (2pt);
\node[fill=white,inner sep=1pt,above right] at (0.4000,7.3000) {$q_S$};
\fill[robotframe] (7.5000,1.3000) circle (2pt);
\node[fill=white,inner sep=1pt,above right] at (7.5000,1.3000) {$q_G$};
\begin{scope}[xshift=9.5cm]
\draw[black!20] (0,0) rectangle (8,8);
\draw[black!50] (5.0000,1.0000) -- (2.0000,2.0000);
\draw[black!50] (7.0000,1.0000) -- (5.0000,1.0000);
\draw[black!50] (7.0000,3.0000) -- (7.0000,1.0000);
\draw[black!50] (1.0000,5.0000) -- (2.0000,2.0000);
\draw[black!50] (3.0000,5.0000) -- (2.0000,2.0000);
\draw[black!50] (3.0000,5.0000) -- (1.0000,5.0000);
\draw[black!50] (1.0000,7.0000) -- (1.0000,5.0000);
\draw[black!50] (3.0000,7.0000) -- (3.0000,5.0000);
\draw[black!50] (3.0000,7.0000) -- (1.0000,7.0000);
\draw[black!50] (7.0000,5.0000) -- (7.0000,3.0000);
\draw[black!50] (5.0000,7.0000) -- (3.0000,7.0000);
\draw[black!50] (7.0000,7.0000) -- (7.0000,5.0000);
\draw[black!50] (7.0000,7.0000) -- (5.0000,7.0000);
\draw[robotframe,very thick] (1.0000,7.0000) -- (1.0000,5.0000) -- (2.0000,2.0000) -- (5.0000,1.0000) -- (7.0000,1.0000);
\draw[fill=white] (2.0000,2.0000) circle (3pt);
\draw[fill=white] (5.0000,1.0000) circle (3pt);
\draw[fill=white] (7.0000,1.0000) circle (3pt);
\draw[fill=black!40] (7.0000,3.0000) circle (3pt);
\draw[fill=black!40] (1.0000,5.0000) circle (3pt);
\draw[fill=black!40] (3.0000,5.0000) circle (3pt);
\draw[fill=white] (1.0000,7.0000) circle (3pt);
\draw[fill=black!40] (3.0000,7.0000) circle (3pt);
\draw[fill=black!40] (7.0000,5.0000) circle (3pt);
\draw[fill=black!40] (5.0000,7.0000) circle (3pt);
\draw[fill=black!40] (7.0000,7.0000) circle (3pt);
\node[above left,fill=white,inner sep=1pt] at (1.0000,7.0000) {$C_S$};
\node[above left,fill=white,inner sep=1pt] at (7.0000,1.0000) {$C_G$};
\end{scope}
\end{scope}
\begin{scope}[yshift=-18.8cm]
\node[anchor=west,font=\small\bfseries] at (0,8.6) {Iteration 3: 34 cells};
\fill[white] (0.0000,0.0000) rectangle (4.0000,4.0000);
\fill[white] (4.0000,0.0000) rectangle (6.0000,2.0000);
\fill[white] (6.0000,0.0000) rectangle (8.0000,2.0000);
\fill[black!65] (4.0000,2.0000) rectangle (6.0000,4.0000);
\fill[black!12] (6.0000,2.0000) rectangle (7.0000,3.0000);
\fill[white] (7.0000,2.0000) rectangle (8.0000,3.0000);
\fill[black!12] (6.0000,3.0000) rectangle (7.0000,4.0000);
\fill[white] (7.0000,3.0000) rectangle (8.0000,4.0000);
\fill[white] (0.0000,4.0000) rectangle (1.0000,5.0000);
\fill[white] (1.0000,4.0000) rectangle (2.0000,5.0000);
\fill[white] (0.0000,5.0000) rectangle (1.0000,6.0000);
\fill[black!12] (1.0000,5.0000) rectangle (2.0000,6.0000);
\fill[white] (2.0000,4.0000) rectangle (3.0000,5.0000);
\fill[white] (3.0000,4.0000) rectangle (4.0000,5.0000);
\fill[black!12] (2.0000,5.0000) rectangle (3.0000,6.0000);
\fill[black!12] (3.0000,5.0000) rectangle (4.0000,6.0000);
\fill[white] (0.0000,6.0000) rectangle (2.0000,8.0000);
\fill[black!12] (2.0000,6.0000) rectangle (3.0000,7.0000);
\fill[black!12] (3.0000,6.0000) rectangle (4.0000,7.0000);
\fill[white] (2.0000,7.0000) rectangle (3.0000,8.0000);
\fill[white] (3.0000,7.0000) rectangle (4.0000,8.0000);
\fill[black!65] (4.0000,4.0000) rectangle (6.0000,6.0000);
\fill[black!12] (6.0000,4.0000) rectangle (7.0000,5.0000);
\fill[white] (7.0000,4.0000) rectangle (8.0000,5.0000);
\fill[black!12] (6.0000,5.0000) rectangle (7.0000,6.0000);
\fill[white] (7.0000,5.0000) rectangle (8.0000,6.0000);
\fill[white] (4.0000,6.0000) rectangle (5.0000,7.0000);
\fill[black!12] (5.0000,6.0000) rectangle (6.0000,7.0000);
\fill[white] (4.0000,7.0000) rectangle (5.0000,8.0000);
\fill[white] (5.0000,7.0000) rectangle (6.0000,8.0000);
\fill[black!12] (6.0000,6.0000) rectangle (7.0000,7.0000);
\fill[white] (7.0000,6.0000) rectangle (8.0000,7.0000);
\fill[white] (6.0000,7.0000) rectangle (7.0000,8.0000);
\fill[white] (7.0000,7.0000) rectangle (8.0000,8.0000);
\fill[black!85] (1.8000,5.2000) -- (2.7000,7.0000) -- (3.7000,5.2000) -- cycle;
\fill[black!85] (4.0000,2.0000) -- (6.0000,2.0000) -- (6.0000,2.5000) -- (6.8000,2.5000) -- (6.8000,6.7000) -- (5.0000,6.7000) -- (5.0000,6.0000) -- (4.0000,6.0000) -- cycle;
\draw[black!45,thin] (0.0000,0.0000) rectangle (4.0000,4.0000);
\draw[black!45,thin] (4.0000,0.0000) rectangle (6.0000,2.0000);
\draw[black!45,thin] (6.0000,0.0000) rectangle (8.0000,2.0000);
\draw[black!45,thin] (4.0000,2.0000) rectangle (6.0000,4.0000);
\draw[black!45,thin] (6.0000,2.0000) rectangle (7.0000,3.0000);
\draw[black!45,thin] (7.0000,2.0000) rectangle (8.0000,3.0000);
\draw[black!45,thin] (6.0000,3.0000) rectangle (7.0000,4.0000);
\draw[black!45,thin] (7.0000,3.0000) rectangle (8.0000,4.0000);
\draw[black!45,thin] (0.0000,4.0000) rectangle (1.0000,5.0000);
\draw[black!45,thin] (1.0000,4.0000) rectangle (2.0000,5.0000);
\draw[black!45,thin] (0.0000,5.0000) rectangle (1.0000,6.0000);
\draw[black!45,thin] (1.0000,5.0000) rectangle (2.0000,6.0000);
\draw[black!45,thin] (2.0000,4.0000) rectangle (3.0000,5.0000);
\draw[black!45,thin] (3.0000,4.0000) rectangle (4.0000,5.0000);
\draw[black!45,thin] (2.0000,5.0000) rectangle (3.0000,6.0000);
\draw[black!45,thin] (3.0000,5.0000) rectangle (4.0000,6.0000);
\draw[black!45,thin] (0.0000,6.0000) rectangle (2.0000,8.0000);
\draw[black!45,thin] (2.0000,6.0000) rectangle (3.0000,7.0000);
\draw[black!45,thin] (3.0000,6.0000) rectangle (4.0000,7.0000);
\draw[black!45,thin] (2.0000,7.0000) rectangle (3.0000,8.0000);
\draw[black!45,thin] (3.0000,7.0000) rectangle (4.0000,8.0000);
\draw[black!45,thin] (4.0000,4.0000) rectangle (6.0000,6.0000);
\draw[black!45,thin] (6.0000,4.0000) rectangle (7.0000,5.0000);
\draw[black!45,thin] (7.0000,4.0000) rectangle (8.0000,5.0000);
\draw[black!45,thin] (6.0000,5.0000) rectangle (7.0000,6.0000);
\draw[black!45,thin] (7.0000,5.0000) rectangle (8.0000,6.0000);
\draw[black!45,thin] (4.0000,6.0000) rectangle (5.0000,7.0000);
\draw[black!45,thin] (5.0000,6.0000) rectangle (6.0000,7.0000);
\draw[black!45,thin] (4.0000,7.0000) rectangle (5.0000,8.0000);
\draw[black!45,thin] (5.0000,7.0000) rectangle (6.0000,8.0000);
\draw[black!45,thin] (6.0000,6.0000) rectangle (7.0000,7.0000);
\draw[black!45,thin] (7.0000,6.0000) rectangle (8.0000,7.0000);
\draw[black!45,thin] (6.0000,7.0000) rectangle (7.0000,8.0000);
\draw[black!45,thin] (7.0000,7.0000) rectangle (8.0000,8.0000);
\draw[thick] (0,0) rectangle (8,8);
\draw[robotframe,very thick] (0.4000,7.3000) -- (1.0000,7.0000) -- (0.5000,6.0000) -- (0.5000,5.5000) -- (0.5000,5.0000) -- (0.5000,4.5000) -- (0.5000,4.0000) -- (2.0000,2.0000) -- (4.0000,1.0000) -- (5.0000,1.0000) -- (6.0000,1.0000) -- (7.0000,1.0000) -- (7.5000,1.3000);
\fill[robotframe] (0.4000,7.3000) circle (2pt);
\node[fill=white,inner sep=1pt,above right] at (0.4000,7.3000) {$q_S$};
\fill[robotframe] (7.5000,1.3000) circle (2pt);
\node[fill=white,inner sep=1pt,above right] at (7.5000,1.3000) {$q_G$};
\begin{scope}[xshift=9.5cm]
\draw[black!20] (0,0) rectangle (8,8);
\draw[black!50] (5.0000,1.0000) -- (2.0000,2.0000);
\draw[black!50] (7.0000,1.0000) -- (5.0000,1.0000);
\draw[black!50] (6.5000,2.5000) -- (7.0000,1.0000);
\draw[black!50] (7.5000,2.5000) -- (7.0000,1.0000);
\draw[black!50] (7.5000,2.5000) -- (6.5000,2.5000);
\draw[black!50] (6.5000,3.5000) -- (6.5000,2.5000);
\draw[black!50] (7.5000,3.5000) -- (7.5000,2.5000);
\draw[black!50] (7.5000,3.5000) -- (6.5000,3.5000);
\draw[black!50] (0.5000,4.5000) -- (2.0000,2.0000);
\draw[black!50] (1.5000,4.5000) -- (2.0000,2.0000);
\draw[black!50] (1.5000,4.5000) -- (0.5000,4.5000);
\draw[black!50] (0.5000,5.5000) -- (0.5000,4.5000);
\draw[black!50] (1.5000,5.5000) -- (1.5000,4.5000);
\draw[black!50] (1.5000,5.5000) -- (0.5000,5.5000);
\draw[black!50] (2.5000,4.5000) -- (2.0000,2.0000);
\draw[black!50] (2.5000,4.5000) -- (1.5000,4.5000);
\draw[black!50] (3.5000,4.5000) -- (2.0000,2.0000);
\draw[black!50] (3.5000,4.5000) -- (2.5000,4.5000);
\draw[black!50] (2.5000,5.5000) -- (1.5000,5.5000);
\draw[black!50] (2.5000,5.5000) -- (2.5000,4.5000);
\draw[black!50] (3.5000,5.5000) -- (3.5000,4.5000);
\draw[black!50] (3.5000,5.5000) -- (2.5000,5.5000);
\draw[black!50] (1.0000,7.0000) -- (0.5000,5.5000);
\draw[black!50] (1.0000,7.0000) -- (1.5000,5.5000);
\draw[black!50] (2.5000,6.5000) -- (2.5000,5.5000);
\draw[black!50] (2.5000,6.5000) -- (1.0000,7.0000);
\draw[black!50] (3.5000,6.5000) -- (3.5000,5.5000);
\draw[black!50] (3.5000,6.5000) -- (2.5000,6.5000);
\draw[black!50] (2.5000,7.5000) -- (1.0000,7.0000);
\draw[black!50] (2.5000,7.5000) -- (2.5000,6.5000);
\draw[black!50] (3.5000,7.5000) -- (3.5000,6.5000);
\draw[black!50] (3.5000,7.5000) -- (2.5000,7.5000);
\draw[black!50] (6.5000,4.5000) -- (6.5000,3.5000);
\draw[black!50] (7.5000,4.5000) -- (7.5000,3.5000);
\draw[black!50] (7.5000,4.5000) -- (6.5000,4.5000);
\draw[black!50] (6.5000,5.5000) -- (6.5000,4.5000);
\draw[black!50] (7.5000,5.5000) -- (7.5000,4.5000);
\draw[black!50] (7.5000,5.5000) -- (6.5000,5.5000);
\draw[black!50] (4.5000,6.5000) -- (3.5000,6.5000);
\draw[black!50] (5.5000,6.5000) -- (4.5000,6.5000);
\draw[black!50] (4.5000,7.5000) -- (3.5000,7.5000);
\draw[black!50] (4.5000,7.5000) -- (4.5000,6.5000);
\draw[black!50] (5.5000,7.5000) -- (5.5000,6.5000);
\draw[black!50] (5.5000,7.5000) -- (4.5000,7.5000);
\draw[black!50] (6.5000,6.5000) -- (6.5000,5.5000);
\draw[black!50] (6.5000,6.5000) -- (5.5000,6.5000);
\draw[black!50] (7.5000,6.5000) -- (7.5000,5.5000);
\draw[black!50] (7.5000,6.5000) -- (6.5000,6.5000);
\draw[black!50] (6.5000,7.5000) -- (5.5000,7.5000);
\draw[black!50] (6.5000,7.5000) -- (6.5000,6.5000);
\draw[black!50] (7.5000,7.5000) -- (7.5000,6.5000);
\draw[black!50] (7.5000,7.5000) -- (6.5000,7.5000);
\draw[robotframe,very thick] (1.0000,7.0000) -- (0.5000,5.5000) -- (0.5000,4.5000) -- (2.0000,2.0000) -- (5.0000,1.0000) -- (7.0000,1.0000);
\draw[fill=white] (2.0000,2.0000) circle (3pt);
\draw[fill=white] (5.0000,1.0000) circle (3pt);
\draw[fill=white] (7.0000,1.0000) circle (3pt);
\draw[fill=black!40] (6.5000,2.5000) circle (3pt);
\draw[fill=white] (7.5000,2.5000) circle (3pt);
\draw[fill=black!40] (6.5000,3.5000) circle (3pt);
\draw[fill=white] (7.5000,3.5000) circle (3pt);
\draw[fill=white] (0.5000,4.5000) circle (3pt);
\draw[fill=white] (1.5000,4.5000) circle (3pt);
\draw[fill=white] (0.5000,5.5000) circle (3pt);
\draw[fill=black!40] (1.5000,5.5000) circle (3pt);
\draw[fill=white] (2.5000,4.5000) circle (3pt);
\draw[fill=white] (3.5000,4.5000) circle (3pt);
\draw[fill=black!40] (2.5000,5.5000) circle (3pt);
\draw[fill=black!40] (3.5000,5.5000) circle (3pt);
\draw[fill=white] (1.0000,7.0000) circle (3pt);
\draw[fill=black!40] (2.5000,6.5000) circle (3pt);
\draw[fill=black!40] (3.5000,6.5000) circle (3pt);
\draw[fill=white] (2.5000,7.5000) circle (3pt);
\draw[fill=white] (3.5000,7.5000) circle (3pt);
\draw[fill=black!40] (6.5000,4.5000) circle (3pt);
\draw[fill=white] (7.5000,4.5000) circle (3pt);
\draw[fill=black!40] (6.5000,5.5000) circle (3pt);
\draw[fill=white] (7.5000,5.5000) circle (3pt);
\draw[fill=white] (4.5000,6.5000) circle (3pt);
\draw[fill=black!40] (5.5000,6.5000) circle (3pt);
\draw[fill=white] (4.5000,7.5000) circle (3pt);
\draw[fill=white] (5.5000,7.5000) circle (3pt);
\draw[fill=black!40] (6.5000,6.5000) circle (3pt);
\draw[fill=white] (7.5000,6.5000) circle (3pt);
\draw[fill=white] (6.5000,7.5000) circle (3pt);
\draw[fill=white] (7.5000,7.5000) circle (3pt);
\node[above left,fill=white,inner sep=1pt] at (1.0000,7.0000) {$C_S$};
\node[above left,fill=white,inner sep=1pt] at (7.0000,1.0000) {$C_G$};
\end{scope}
\end{scope}
\end{tikzpicture}

\caption{Approximate decomposition and its connectivity graphs, using the
same obstacles at all three iterations. White cells/nodes are empty,
gray ones are mixed, and black cells are occupied and excluded from the
graphs. Red graph routes in iterations 1 and 2 are provisional because
they visit mixed cells. The third route uses only empty cells; the map
shows a validated geometric realization through free portals and cell
interiors. Graph edges describe adjacency, not executable straight paths
between cell centers when mixed cells are involved.}
\label{fig:l13-quadtree}
\end{figure}

\subsection{From cells to a path}
For adjacent empty cells, the free part of their shared side is a
\emph{portal}. Its midpoint is a convenient crossing point when the
whole portal is free. Join consecutive crossings through the cell interior,
inserting an interior waypoint if a direct connection would follow a
forbidden boundary. Square cells are convex, so interior waypoints make
this construction particularly simple. Include connections from $q_S$
and to $q_G$.

Refinement requires a stopping rule, such as a maximum depth or minimum
cell width. Failure to obtain a certified route at that resolution does
not generally prove that no continuous path exists. A narrow passage can
remain hidden in mixed cells. Conversely, a graph disconnected even when
all mixed cells are allowed can certify separation if that graph includes
every possible transition under the chosen domain assumptions.

A shortest route on the cell graph minimizes its assigned graph cost.
It need not be the shortest geometric path. Refining only mixed cells
does not by itself establish convergence to the continuous shortest path.
After constructing a feasible path, one can try collision-checked
\emph{shortcuts}: replace a subpath by a straight segment only when the
entire replacement is free. This can reduce length without certifying
global optimality.

\section{Probabilistic roadmaps}
\label{sec:l13-prm}
Cell decompositions construct a representation by partitioning space.
A \emph{probabilistic roadmap} (PRM) instead samples configurations and
connects nearby samples with feasible local paths. Different samples can
produce different graphs, so adding samples can reveal connections that
were absent in an earlier attempt.

For a point moving in a planar free region, a local path can be a straight
segment. A simple construction is:
\begin{enumerate}
\item Sample a batch of $N$ configurations in the domain.
\item Reject configurations in collision. Retain the free samples as nodes.
\item For pairs closer than a connection radius $d$, test the complete
straight segment between them. Add an edge only if that segment is free.
Assign a cost, such as its Euclidean length.
\item Insert $q_S$ and $q_G$ and connect them by the same local test.
Search the graph using Dijkstra or another suitable graph-search method.
\item If no route exists in the current graph, add more samples and
connections, subject to a computation budget, and search again.
\end{enumerate}
In a fixed environment the roadmap can be reused for different start--goal
queries. Only the query configurations and their connections need to be
inserted again, provided the robot model and collision geometry are unchanged.

\begin{figure}[htbp]
\centering% Generated by tools/generate_sampling_figures.py
\begin{tikzpicture}[scale=1.35,line cap=round,line join=round,
 every node/.style={font=\footnotesize},>=Stealth]
\draw[thick] (0.0000,0.0000) -- (10.0000,0.0000) -- (9.3000,3.7000) -- (3.0000,3.3000) -- cycle;
\fill[black!85] (2.5000,1.0000) -- (4.0000,2.5000) -- (4.5000,0.8000) -- cycle;
\fill[black!85] (6.0000,1.2000) -- (7.8000,1.2000) -- (8.0000,2.7000) -- (6.1000,2.7000) -- cycle;
\draw[black!40] (4.1000,0.4500) -- (2.1000,0.6500);
\draw[black!40] (6.0000,3.0000) -- (4.9000,1.8000);
\draw[black!40] (8.3000,2.1000) -- (8.5000,0.6000);
\draw[black!40] (2.0000,1.5000) -- (3.0000,2.9500);
\draw[black!40] (2.0000,1.5000) -- (2.1000,0.6500);
\draw[black!40] (4.6000,2.9000) -- (3.0000,2.9500);
\draw[black!40] (4.6000,2.9000) -- (4.9000,1.8000);
\draw[black!40] (4.6000,2.9000) -- (6.0000,3.0000);
\draw[black!40] (5.7000,0.5000) -- (4.1000,0.4500);
\draw[black!40] (5.7000,0.5000) -- (4.9000,1.8000);
\draw[black!40] (7.0000,0.5000) -- (8.5000,0.6000);
\draw[black!40] (7.0000,0.5000) -- (5.7000,0.5000);
\draw[black!40] (8.7000,3.1000) -- (8.3000,2.1000);
\draw[black!40] (2.8000,2.0000) -- (3.0000,2.9500);
\draw[black!40] (2.8000,2.0000) -- (2.1000,0.6500);
\draw[black!40] (2.8000,2.0000) -- (2.0000,1.5000);
\draw[black!40] (2.8000,2.0000) -- (4.6000,2.9000);
\draw[black!40] (3.0000,0.4000) -- (2.1000,0.6500);
\draw[black!40] (3.0000,0.4000) -- (4.1000,0.4500);
\draw[black!40] (3.0000,0.4000) -- (2.0000,1.5000);
\draw[black!40] (5.8000,1.6000) -- (4.1000,0.4500);
\draw[black!40] (5.8000,1.6000) -- (4.9000,1.8000);
\draw[black!40] (5.8000,1.6000) -- (6.0000,3.0000);
\draw[black!40] (5.8000,1.6000) -- (4.6000,2.9000);
\draw[black!40] (5.8000,1.6000) -- (5.7000,0.5000);
\draw[black!40] (8.6000,1.3000) -- (8.5000,0.6000);
\draw[black!40] (8.6000,1.3000) -- (8.3000,2.1000);
\draw[black!40] (8.6000,1.3000) -- (7.0000,0.5000);
\draw[black!40] (8.6000,1.3000) -- (8.7000,3.1000);
\draw[black!40] (7.4000,3.2000) -- (6.0000,3.0000);
\draw[black!40] (7.4000,3.2000) -- (8.7000,3.1000);
\draw[red!75!black,dashed,thick] (2.8000,2.0000) -- (4.9000,1.8000);
\draw[robotframe,very thick] (3.0000,2.9500) -- (4.6000,2.9000) -- (4.9000,1.8000) -- (5.7000,0.5000) -- (7.0000,0.5000) -- (8.5000,0.6000);
\draw[black,fill=white,thick] (3.0000,2.9500) circle (2.5pt);
\draw[black,fill=white,thick] (8.5000,0.6000) circle (2.5pt);
\draw[sensorframe,fill=white,thick] (2.1000,0.6500) circle (2.5pt);
\draw[sensorframe,fill=white,thick] (4.1000,0.4500) circle (2.5pt);
\draw[sensorframe,fill=white,thick] (4.9000,1.8000) circle (2.5pt);
\draw[sensorframe,fill=white,thick] (6.0000,3.0000) circle (2.5pt);
\draw[sensorframe,fill=white,thick] (8.3000,2.1000) circle (2.5pt);
\draw[orange!90!black,fill=white,thick] (2.0000,1.5000) circle (2.5pt);
\draw[orange!90!black,fill=white,thick] (4.6000,2.9000) circle (2.5pt);
\draw[orange!90!black,fill=white,thick] (5.7000,0.5000) circle (2.5pt);
\draw[orange!90!black,fill=white,thick] (7.0000,0.5000) circle (2.5pt);
\draw[orange!90!black,fill=white,thick] (8.7000,3.1000) circle (2.5pt);
\draw[green!50!black,fill=white,thick] (2.8000,2.0000) circle (2.5pt);
\draw[green!50!black,fill=white,thick] (3.0000,0.4000) circle (2.5pt);
\draw[green!50!black,fill=white,thick] (5.8000,1.6000) circle (2.5pt);
\draw[green!50!black,fill=white,thick] (8.6000,1.3000) circle (2.5pt);
\draw[green!50!black,fill=white,thick] (7.4000,3.2000) circle (2.5pt);
\node[red,font=\large] at (3.7000,1.5000) {$\times$};
\node[red,font=\large] at (6.8000,1.8000) {$\times$};
\node[above] at (3.0000,2.9500) {$q_S$};
\node[below] at (8.5000,0.6000) {$q_G$};
\node[anchor=west] at (0,-.5) {Samples: \textcolor{sensorframe}{batch 1}, \textcolor{orange!90!black}{batch 2}, \textcolor{green!50!black}{batch 3}};
\end{tikzpicture}

\caption{A sampled roadmap in a polygonal domain. Sample colors distinguish
three batches; red crosses mark rejected configurations inside obstacles.
The dashed red candidate edge has free endpoints but crosses the triangular
obstacle, so it is rejected. The highlighted route consists entirely of
validated edges. Coordinates are illustrative.}
\label{fig:l13-prm}
\end{figure}

\subsection{Checking nodes is not enough}
Two free configurations can lie on opposite sides of an obstacle. Thus
node validity and edge validity are separate tests. If an attempted edge
is blocked, remove that edge, not its otherwise valid endpoints: those
nodes may still connect to other neighbors. The same distinction applies
when an edge is a curved robot motion rather than a straight line.

The radius $d$ controls the candidate connections. A small radius reduces
the number of attempted edges but can leave the graph disconnected.
A larger radius offers more connections and requires more collision tests.
For configurations involving angles or mixed units, the distance used to
select neighbors must be defined appropriately; it is not automatically
ordinary distance between unscaled coordinate entries.

\subsection{What randomness does and does not guarantee}
An unsuccessful finite run means that the current samples and connections
did not produce a route. It does not establish that the planning problem
is impossible. Biasing samples toward the goal or a likely corridor can
help, but excluding other regions can remove the only feasible detour.

Two different properties must be distinguished:
\begin{itemize}
\item \textbf{Probabilistic completeness:} for a feasible problem satisfying
the method's assumptions, the probability of finding a solution approaches
one as the number of samples grows.
\item \textbf{Asymptotic optimality:} under the relevant assumptions, the
best solution cost converges to the optimal cost as sampling continues.
\end{itemize}
Neither property means that a finite run must succeed or return the
continuous optimum. Dijkstra finds the optimum on the constructed graph;
that alone proves neither property for the sampling procedure. Guarantees
depend on the sampling distribution, local planner, connection rule, and
geometry. PRM* is a roadmap variant designed for asymptotic optimality.

\section{Growing search trees}
\label{sec:l13-trees}
A roadmap may connect many pairs of samples. A tree instead gives each
new node one parent and grows outward from an initial configuration.
Once a node has a valid connection to the goal, its parent chain recovers
a path to the root. There is no need to run Dijkstra merely to recover
the unique path already stored in that tree.

\subsection{A local branching construction}
A useful geometric construction starts at $q_S$, proposes a few points
at distance $d$ from it, and retains the collision-free connecting segments.
Repeat from the new nodes, drawing successive generations in different
colors. Reject candidates whose motions hit an obstacle or leave the domain.
An unsuccessful extension does not necessarily make its parent useless:
another direction from the same parent may be valid.

If every accepted edge has length $d$, a node at depth $k$ has path
length $kd$. Expanding by increasing depth then agrees with increasing
path cost. If lengths differ, use accumulated edge costs instead. A
priority based on cost plus a suitable heuristic can guide this kind of
search, but it defines a different expansion policy from standard RRT.
Figure~\ref{fig:l13-local-tree} shows the local construction.

\begin{figure}[htbp]
\centering% Generated by tools/generate_sampling_figures.py
\begin{tikzpicture}[scale=1.35,line cap=round,line join=round,
 every node/.style={font=\footnotesize},>=Stealth]
\draw[thick] (0,0) rectangle (9,4.5);
\fill[black!85] (1.4000,2.1000) -- (2.0000,3.4000) -- (2.7000,2.1000) -- cycle;
\fill[black!85] (5.0000,1.1000) -- (6.2000,3.4000) -- (7.5000,1.1000) -- cycle;
\draw[robotframe,thick] (0.4500,4.0000) -- (0.7578,3.1543);
\fill[robotframe] (0.7578,3.1543) circle (1.5pt);
\draw[robotframe,thick] (0.4500,4.0000) -- (1.3363,3.8437);
\fill[robotframe] (1.3363,3.8437) circle (1.5pt);
\draw[red!65,dashed] (0.4500,4.0000) -- (0.7578,4.8457);
\node[red!65] at (0.7578,4.8457) {$\times$};
\draw[red!65,dashed] (0.4500,4.0000) -- (-0.2394,4.5785);
\node[red!65] at (-0.2394,4.5785) {$\times$};
\draw[red!65,dashed] (0.4500,4.0000) -- (-0.3957,3.6922);
\node[red!65] at (-0.3957,3.6922) {$\times$};
\draw[green!50!black,thick] (0.7578,3.1543) -- (0.6856,2.2572);
\fill[green!50!black] (0.6856,2.2572) circle (1.5pt);
\draw[green!50!black,thick] (0.7578,3.1543) -- (1.4525,2.5821);
\fill[green!50!black] (1.4525,2.5821) circle (1.5pt);
\draw[green!50!black,thick] (1.3363,3.8437) -- (2.1403,3.4392);
\fill[green!50!black] (2.1403,3.4392) circle (1.5pt);
\draw[green!50!black,thick] (1.3363,3.8437) -- (2.2191,4.0187);
\fill[green!50!black] (2.2191,4.0187) circle (1.5pt);
\draw[sensorframe,thick] (0.6856,2.2572) -- (0.3079,1.4403);
\fill[sensorframe] (0.3079,1.4403) circle (1.5pt);
\draw[sensorframe,thick] (0.6856,2.2572) -- (1.0527,1.4354);
\fill[sensorframe] (1.0527,1.4354) circle (1.5pt);
\draw[red!65,dashed] (1.4525,2.5821) -- (1.8809,1.7906);
\node[red!65] at (1.8809,1.7906) {$\times$};
\draw[red!65,dashed] (1.4525,2.5821) -- (2.2899,2.2523);
\node[red!65] at (2.2899,2.2523) {$\times$};
\draw[sensorframe,thick] (2.1403,3.4392) -- (2.5872,2.6580);
\fill[sensorframe] (2.5872,2.6580) circle (1.5pt);
\draw[sensorframe,thick] (2.1403,3.4392) -- (3.0383,3.3796);
\fill[sensorframe] (3.0383,3.3796) circle (1.5pt);
\draw[sensorframe,thick] (2.2191,4.0187) -- (3.1072,3.8725);
\fill[sensorframe] (3.1072,3.8725) circle (1.5pt);
\draw[sensorframe,thick] (2.2191,4.0187) -- (2.9844,4.4924);
\fill[sensorframe] (2.9844,4.4924) circle (1.5pt);
\draw[red!65,dashed] (0.3079,1.4403) -- (-0.3794,0.8592);
\node[red!65] at (-0.3794,0.8592) {$\times$};
\draw[violet,thick] (0.3079,1.4403) -- (0.2592,0.5416);
\fill[violet] (0.2592,0.5416) circle (1.5pt);
\draw[violet,thick] (1.0527,1.4354) -- (1.1359,0.5393);
\fill[violet] (1.1359,0.5393) circle (1.5pt);
\draw[violet,thick] (1.0527,1.4354) -- (1.8082,0.9464);
\fill[violet] (1.8082,0.9464) circle (1.5pt);
\draw[red!65,dashed] (2.5872,2.6580) -- (2.6343,1.7592);
\node[red!65] at (2.6343,1.7592) {$\times$};
\draw[violet,thick] (2.5872,2.6580) -- (3.3944,2.2599);
\fill[violet] (3.3944,2.2599) circle (1.5pt);
\draw[violet,thick] (3.0383,3.3796) -- (3.7601,2.8420);
\fill[violet] (3.7601,2.8420) circle (1.5pt);
\draw[violet,thick] (3.0383,3.3796) -- (3.8414,3.7858);
\fill[violet] (3.8414,3.7858) circle (1.5pt);
\draw[violet,thick] (3.1072,3.8725) -- (3.9048,3.4556);
\fill[violet] (3.9048,3.4556) circle (1.5pt);
\draw[violet,thick] (3.1072,3.8725) -- (3.9952,4.0189);
\fill[violet] (3.9952,4.0189) circle (1.5pt);
\draw[red!65,dashed] (2.9844,4.4924) -- (3.8695,4.6557);
\node[red!65] at (3.8695,4.6557) {$\times$};
\draw[red!65,dashed] (2.9844,4.4924) -- (3.4638,5.2541);
\node[red!65] at (3.4638,5.2541) {$\times$};
\draw[red!65,dashed] (0.2592,0.5416) -- (-0.0473,-0.3046);
\node[red!65] at (-0.0473,-0.3046) {$\times$};
\draw[red!65,dashed] (0.2592,0.5416) -- (0.5175,-0.3205);
\node[red!65] at (0.5175,-0.3205) {$\times$};
\draw[red!65,dashed] (1.1359,0.5393) -- (0.8432,-0.3118);
\node[red!65] at (0.8432,-0.3118) {$\times$};
\draw[red!65,dashed] (1.1359,0.5393) -- (1.6799,-0.1777);
\node[red!65] at (1.6799,-0.1777) {$\times$};
\draw[orange!85!black,thick] (1.8082,0.9464) -- (2.3190,0.2054);
\fill[orange!85!black] (2.3190,0.2054) circle (1.5pt);
\draw[orange!85!black,thick] (1.8082,0.9464) -- (2.6906,0.7693);
\fill[orange!85!black] (2.6906,0.7693) circle (1.5pt);
\draw[orange!85!black,thick] (3.3944,2.2599) -- (3.8608,1.4902);
\fill[orange!85!black] (3.8608,1.4902) circle (1.5pt);
\draw[orange!85!black,thick] (3.3944,2.2599) -- (4.2883,2.3643);
\fill[orange!85!black] (4.2883,2.3643) circle (1.5pt);
\draw[orange!85!black,thick] (3.7601,2.8420) -- (4.1807,2.0463);
\fill[orange!85!black] (4.1807,2.0463) circle (1.5pt);
\draw[orange!85!black,thick] (3.7601,2.8420) -- (4.6591,2.7990);
\fill[orange!85!black] (4.6591,2.7990) circle (1.5pt);
\draw[orange!85!black,thick] (3.8414,3.7858) -- (4.7385,3.7133);
\fill[orange!85!black] (4.7385,3.7133) circle (1.5pt);
\draw[red!65,dashed] (3.8414,3.7858) -- (4.3256,4.5444);
\node[red!65] at (4.3256,4.5444) {$\times$};
\draw[orange!85!black,thick] (3.9048,3.4556) -- (4.5191,2.7978);
\fill[orange!85!black] (4.5191,2.7978) circle (1.5pt);
\draw[orange!85!black,thick] (3.9048,3.4556) -- (4.8014,3.5334);
\fill[orange!85!black] (4.8014,3.5334) circle (1.5pt);
\draw[orange!85!black,thick] (3.9952,4.0189) -- (4.8154,3.6484);
\fill[orange!85!black] (4.8154,3.6484) circle (1.5pt);
\draw[orange!85!black,thick] (3.9952,4.0189) -- (4.7579,4.4967);
\fill[orange!85!black] (4.7579,4.4967) circle (1.5pt);
\fill (0.4500,4.0000) circle (2pt) node[below right] {$q_S$};
\fill (4.4,.65) circle (2pt) node[right] {$q_G$};
\node[anchor=west] at (0,-.7) {Generations: \textcolor{robotframe}{1}, \textcolor{green!50!black}{2}, \textcolor{sensorframe}{3}, \textcolor{violet}{4}, \textcolor{orange!85!black}{5}};
\end{tikzpicture}

\caption{Local branching from $q_S$ with equal-length candidate steps.
Colors identify successive generations. Dashed red extensions and crosses
mark rejected candidates; the existing free parent nodes remain valid.
The tree has not yet reached $q_G$. This generation-by-generation
construction illustrates tree growth; it does not use the standard RRT
node-selection rule.}
\label{fig:l13-local-tree}
\end{figure}

\subsection{Rapidly exploring random trees}
An \emph{RRT}, or \emph{rapidly exploring random tree}, selects where to
expand using a sample from the configuration space. For a planar geometric
problem, one iteration is:
\begin{enumerate}
\item Sample a target configuration $q_{\mathrm{rand}}$.
\item Find the tree vertex $q_{\mathrm{near}}$ closest to that target.
\item Steer from $q_{\mathrm{near}}$ toward $q_{\mathrm{rand}}$ by at most
a chosen step length $d$, obtaining a candidate $q_{\mathrm{new}}$.
\item If the complete extension is free, add $q_{\mathrm{new}}$ with
parent $q_{\mathrm{near}}$; otherwise reject the extension.
\end{enumerate}
In Euclidean coordinates, for a nonzero separation, the steering step is
\[
 q_{\mathrm{new}}=q_{\mathrm{near}}+
 \min\!\left(1,\frac{d}{\|q_{\mathrm{rand}}-q_{\mathrm{near}}\|}\right)
 (q_{\mathrm{rand}}-q_{\mathrm{near}}).
\]
Figure~\ref{fig:l13-rrt-step} separates the random target from the new
vertex actually added. The target need not become a tree node.

\begin{figure}[htbp]
\centering% Generated by tools/generate_sampling_figures.py
\begin{tikzpicture}[scale=1.35,line cap=round,line join=round,
 every node/.style={font=\footnotesize},>=Stealth]
\draw[black!20] (0,0) rectangle (8,4);
\draw[thick] (0.5000,0.4000) -- (1.3000,1.3000);
\draw[thick] (1.3000,1.3000) -- (1.0000,2.2000);
\draw[thick] (1.3000,1.3000) -- (2.5000,1.8000);
\draw[thick] (2.5000,1.8000) -- (3.5000,1.1000);
\draw[thick] (2.5000,1.8000) -- (3.0000,3.0000);
\fill (0.5000,0.4000) circle (2pt);
\fill (1.3000,1.3000) circle (2pt);
\fill (1.0000,2.2000) circle (2pt);
\fill (2.5000,1.8000) circle (2pt);
\fill (3.5000,1.1000) circle (2pt);
\fill (3.0000,3.0000) circle (2pt);
\draw[sensorframe,dashed] (4.7868,1.6515) -- (7.0000,2.6000);
\draw[sensorframe,very thick,->] (3.5000,1.1000) -- (4.7868,1.6515) node[midway,above left] {$d$};
\fill[sensorframe] (4.7868,1.6515) circle (2.5pt) node[below right] {$q_{\mathrm{new}}$};
\draw[sensorframe,thick] (7.0000,2.6000) circle (3pt) node[above] {$q_{\mathrm{rand}}$};
\node[below] at (3.5000,1.1000) {$q_{\mathrm{near}}$};
\node[below right] at (0.5000,0.4000) {$q_S$};
\end{tikzpicture}

\caption{One RRT extension. A global random target selects its nearest
existing tree vertex. Only the short, validated extension to
$q_{\mathrm{new}}$ is inserted; the dashed remainder is not a tree edge.}
\label{fig:l13-rrt-step}
\end{figure}

Each vertex has a region of targets for which it is the nearest vertex:
its Voronoi region relative to the tree vertices. Under uniform sampling,
vertices with larger such regions are selected more frequently. This
helps exploration spread into uncovered space. RRT therefore does not
simply expand every leaf once per generation.

The goal may be represented by a small goal region, or the planner may
attempt a checked connection to $q_G$ from a nearby vertex. Occasionally
sampling the goal itself is another useful bias. Closeness alone is not
a valid final edge: that connection must satisfy collision and motion
constraints. Basic RRT seeks a feasible path and is not generally
asymptotically optimal; RRT* adds a cost-aware parent choice and rewiring.

\subsection{Bidirectional growth}
Two trees can be rooted at $q_S$ and $q_G$ and grown toward one another.
When vertices in different trees can be joined by a valid local path,
concatenate the start-tree path, the connecting motion, and the goal-tree
path. In a reversible geometric problem the goal-tree edges can be
traversed in reverse order.

\begin{figure}[htbp]
\centering% Generated by tools/generate_sampling_figures.py
\begin{tikzpicture}[scale=1.35,line cap=round,line join=round,
 every node/.style={font=\footnotesize},>=Stealth]
\draw[thick] (0,0) rectangle (9,4.8);
\fill[black!85] (2.6000,0.9000) -- (6.1000,0.9000) -- (6.4000,2.7000) -- (3.0000,2.8000) -- cycle;
\draw[sensorframe,thick] (0.6000,3.5000) -- (1.5000,3.7000);
\draw[sensorframe,thick] (1.5000,3.7000) -- (2.4000,4.0000);
\draw[sensorframe,thick] (2.4000,4.0000) -- (3.5000,4.1000);
\draw[sensorframe,thick] (0.6000,3.5000) -- (1.1000,2.4000);
\draw[sensorframe,thick] (1.1000,2.4000) -- (0.5000,1.5000);
\draw[sensorframe,thick] (1.1000,2.4000) -- (1.8000,1.7000);
\draw[sensorframe,thick] (1.8000,1.7000) -- (2.3000,1.2000);
\fill[sensorframe] (0.6000,3.5000) circle (2pt);
\fill[sensorframe] (1.5000,3.7000) circle (2pt);
\fill[sensorframe] (2.4000,4.0000) circle (2pt);
\fill[sensorframe] (3.5000,4.1000) circle (2pt);
\fill[sensorframe] (1.1000,2.4000) circle (2pt);
\fill[sensorframe] (0.5000,1.5000) circle (2pt);
\fill[sensorframe] (1.8000,1.7000) circle (2pt);
\fill[sensorframe] (2.3000,1.2000) circle (2pt);
\draw[green!50!black,thick] (8.5000,3.4000) -- (7.8000,3.9000);
\draw[green!50!black,thick] (7.8000,3.9000) -- (6.8000,4.1000);
\draw[green!50!black,thick] (6.8000,4.1000) -- (5.5000,4.1000);
\draw[green!50!black,thick] (8.5000,3.4000) -- (8.0000,2.4000);
\draw[green!50!black,thick] (8.0000,2.4000) -- (8.5000,1.5000);
\draw[green!50!black,thick] (8.0000,2.4000) -- (7.3000,1.5000);
\draw[green!50!black,thick] (7.3000,1.5000) -- (6.8000,0.5000);
\fill[green!50!black] (8.5000,3.4000) circle (2pt);
\fill[green!50!black] (7.8000,3.9000) circle (2pt);
\fill[green!50!black] (6.8000,4.1000) circle (2pt);
\fill[green!50!black] (5.5000,4.1000) circle (2pt);
\fill[green!50!black] (8.0000,2.4000) circle (2pt);
\fill[green!50!black] (8.5000,1.5000) circle (2pt);
\fill[green!50!black] (7.3000,1.5000) circle (2pt);
\fill[green!50!black] (6.8000,0.5000) circle (2pt);
\draw[robotframe,very thick] (0.6000,3.5000) -- (1.5000,3.7000) -- (2.4000,4.0000) -- (3.5000,4.1000) -- (5.5000,4.1000) -- (6.8000,4.1000) -- (7.8000,3.9000) -- (8.5000,3.4000);
\node[above] at (4.5,4.1) {connection};
\node[below left] at (0.6000,3.5000) {$q_S$};
\node[below right] at (8.5000,3.4000) {$q_G$};
\end{tikzpicture}

\caption{Two geometric trees exploring around an obstacle. A checked
connection above the obstacle joins the trees. The highlighted route
combines their parent chains and the connecting edge.}
\label{fig:l13-bidirectional}
\end{figure}

Growing from both ends can reduce the exploration needed before a
connection is found, although the benefit depends on the environment and
extension policy. It does not guarantee a fixed speedup or an optimal path.
For robots with directional motion constraints, constructing the goal
tree requires more care, as the next section explains.

\section{Motion primitives and nonholonomic robots}
\label{sec:l13-primitives}
A free path for a point in the plane is not automatically an executable
motion for a robot. For a unicycle with $q=(x,y,\theta)$,
\[
 \dot{x}=v\cos\theta,\qquad
 \dot{y}=v\sin\theta,\qquad
 \dot{\theta}=\omega.
\]
It cannot translate instantaneously sideways relative to its heading.
An ideal unicycle can rotate in place by choosing $v=0$, but that rotation
takes time for bounded $\omega$. Other platforms, such as a car with
bounded steering, cannot generally rotate in place.

A \emph{motion primitive} is a short feasible motion generated by a
specified control input. For example, hold $v>0$ fixed over a duration
$\Delta t$ and choose one of
\[
 \omega\in\{-\omega_0,\,0,\,+\omega_0\}.
\]
The resulting motions are a right arc, a straight segment, and a left arc.
For $\omega\ne0$, their endpoints are
\begin{align*}
 x'&=x+\frac{v}{\omega}
       [\sin(\theta+\omega\Delta t)-\sin\theta],\\
 y'&=y-\frac{v}{\omega}
       [\cos(\theta+\omega\Delta t)-\cos\theta],\\
 \theta'&=\theta+\omega\Delta t.
\end{align*}
For $\omega=0$, use $x'=x+v\Delta t\cos\theta$,
$y'=y+v\Delta t\sin\theta$, and $\theta'=\theta$.
These are the same constant-input motions used in odometry, now used
to generate candidate successors during planning.

\subsection{Expanding feasible motions}
From a node, generate its permitted primitives and collision-check each
complete swept motion. Retain valid endpoints together with their headings,
parent pointers, and generating controls. Expanding all three primitives
at every surviving node gives a deterministic motion-primitive search
tree. A sampling-based tree can instead use random targets or controls
to choose extensions. Randomness and kinematic feasibility are separate
design choices.

\begin{figure}[htbp]
\centering% Generated by tools/generate_sampling_figures.py
\begin{tikzpicture}[scale=1.3,line cap=round,line join=round,
 every node/.style={font=\footnotesize},>=Stealth]
\draw[thick] (0,0) rectangle (10,5);
\draw[robotframe!100,->] (1.0000,0.7000) -- (1.0008,0.7475) -- (1.0034,0.7949) -- (1.0076,0.8422) -- (1.0135,0.8894) -- (1.0211,0.9362) -- (1.0303,0.9828) -- (1.0412,1.0291) -- (1.0538,1.0749) -- (1.0679,1.1202) -- (1.0837,1.1650) -- (1.1011,1.2092) -- (1.1200,1.2528) -- (1.1405,1.2957) -- (1.1624,1.3378) -- (1.1859,1.3791) -- (1.2108,1.4195) -- (1.2371,1.4590) -- (1.2649,1.4976) -- (1.2940,1.5351) -- (1.3244,1.5716);
\draw[robotframe!100,->] (1.0000,0.7000) -- (1.0000,0.7475) -- (1.0000,0.7950) -- (1.0000,0.8425) -- (1.0000,0.8900) -- (1.0000,0.9375) -- (1.0000,0.9850) -- (1.0000,1.0325) -- (1.0000,1.0800) -- (1.0000,1.1275) -- (1.0000,1.1750) -- (1.0000,1.2225) -- (1.0000,1.2700) -- (1.0000,1.3175) -- (1.0000,1.3650) -- (1.0000,1.4125) -- (1.0000,1.4600) -- (1.0000,1.5075) -- (1.0000,1.5550) -- (1.0000,1.6025) -- (1.0000,1.6500);
\draw[robotframe!100,->] (1.0000,0.7000) -- (0.9992,0.7475) -- (0.9966,0.7949) -- (0.9924,0.8422) -- (0.9865,0.8894) -- (0.9789,0.9362) -- (0.9697,0.9828) -- (0.9588,1.0291) -- (0.9462,1.0749) -- (0.9321,1.1202) -- (0.9163,1.1650) -- (0.8989,1.2092) -- (0.8800,1.2528) -- (0.8595,1.2957) -- (0.8376,1.3378) -- (0.8141,1.3791) -- (0.7892,1.4195) -- (0.7629,1.4590) -- (0.7351,1.4976) -- (0.7060,1.5351) -- (0.6756,1.5716);
\draw[robotframe!80,->] (1.3244,1.5716) -- (1.3560,1.6070) -- (1.3890,1.6413) -- (1.4231,1.6743) -- (1.4584,1.7061) -- (1.4948,1.7366) -- (1.5322,1.7658) -- (1.5707,1.7937) -- (1.6101,1.8202) -- (1.6505,1.8452) -- (1.6917,1.8688) -- (1.7337,1.8909) -- (1.7765,1.9115) -- (1.8200,1.9306) -- (1.8642,1.9481) -- (1.9090,1.9640) -- (1.9542,1.9783) -- (2.0000,1.9910) -- (2.0462,2.0021) -- (2.0928,2.0115) -- (2.1396,2.0192);
\draw[robotframe!80,->] (1.3244,1.5716) -- (1.3554,1.6076) -- (1.3865,1.6435) -- (1.4175,1.6795) -- (1.4486,1.7154) -- (1.4796,1.7514) -- (1.5107,1.7873) -- (1.5417,1.8232) -- (1.5728,1.8592) -- (1.6038,1.8951) -- (1.6349,1.9311) -- (1.6659,1.9670) -- (1.6970,2.0030) -- (1.7280,2.0389) -- (1.7591,2.0749) -- (1.7901,2.1108) -- (1.8212,2.1468) -- (1.8522,2.1827) -- (1.8833,2.2186) -- (1.9143,2.2546) -- (1.9454,2.2905);
\draw[robotframe!80,->] (1.3244,1.5716) -- (1.3548,1.6081) -- (1.3839,1.6457) -- (1.4116,1.6842) -- (1.4379,1.7238) -- (1.4628,1.7642) -- (1.4863,1.8055) -- (1.5083,1.8476) -- (1.5287,1.8905) -- (1.5476,1.9340) -- (1.5650,1.9783) -- (1.5808,2.0231) -- (1.5949,2.0684) -- (1.6075,2.1142) -- (1.6184,2.1604) -- (1.6276,2.2070) -- (1.6352,2.2539) -- (1.6411,2.3010) -- (1.6453,2.3484) -- (1.6479,2.3958) -- (1.6487,2.4433);
\draw[robotframe!80,->] (1.0000,1.6500) -- (1.0008,1.6975) -- (1.0034,1.7449) -- (1.0076,1.7922) -- (1.0135,1.8394) -- (1.0211,1.8862) -- (1.0303,1.9328) -- (1.0412,1.9791) -- (1.0538,2.0249) -- (1.0679,2.0702) -- (1.0837,2.1150) -- (1.1011,2.1592) -- (1.1200,2.2028) -- (1.1405,2.2457) -- (1.1624,2.2878) -- (1.1859,2.3291) -- (1.2108,2.3695) -- (1.2371,2.4090) -- (1.2649,2.4476) -- (1.2940,2.4851) -- (1.3244,2.5216);
\draw[robotframe!80,->] (1.0000,1.6500) -- (1.0000,1.6975) -- (1.0000,1.7450) -- (1.0000,1.7925) -- (1.0000,1.8400) -- (1.0000,1.8875) -- (1.0000,1.9350) -- (1.0000,1.9825) -- (1.0000,2.0300) -- (1.0000,2.0775) -- (1.0000,2.1250) -- (1.0000,2.1725) -- (1.0000,2.2200) -- (1.0000,2.2675) -- (1.0000,2.3150) -- (1.0000,2.3625) -- (1.0000,2.4100) -- (1.0000,2.4575) -- (1.0000,2.5050) -- (1.0000,2.5525) -- (1.0000,2.6000);
\draw[robotframe!80,->] (1.0000,1.6500) -- (0.9992,1.6975) -- (0.9966,1.7449) -- (0.9924,1.7922) -- (0.9865,1.8394) -- (0.9789,1.8862) -- (0.9697,1.9328) -- (0.9588,1.9791) -- (0.9462,2.0249) -- (0.9321,2.0702) -- (0.9163,2.1150) -- (0.8989,2.1592) -- (0.8800,2.2028) -- (0.8595,2.2457) -- (0.8376,2.2878) -- (0.8141,2.3291) -- (0.7892,2.3695) -- (0.7629,2.4090) -- (0.7351,2.4476) -- (0.7060,2.4851) -- (0.6756,2.5216);
\draw[robotframe!80,->] (0.6756,1.5716) -- (0.6452,1.6081) -- (0.6161,1.6457) -- (0.5884,1.6842) -- (0.5621,1.7238) -- (0.5372,1.7642) -- (0.5137,1.8055) -- (0.4917,1.8476) -- (0.4713,1.8905) -- (0.4524,1.9340) -- (0.4350,1.9783) -- (0.4192,2.0231) -- (0.4051,2.0684) -- (0.3925,2.1142) -- (0.3816,2.1604) -- (0.3724,2.2070) -- (0.3648,2.2539) -- (0.3589,2.3010) -- (0.3547,2.3484) -- (0.3521,2.3958) -- (0.3513,2.4433);
\draw[robotframe!80,->] (0.6756,1.5716) -- (0.6446,1.6076) -- (0.6135,1.6435) -- (0.5825,1.6795) -- (0.5514,1.7154) -- (0.5204,1.7514) -- (0.4893,1.7873) -- (0.4583,1.8232) -- (0.4272,1.8592) -- (0.3962,1.8951) -- (0.3651,1.9311) -- (0.3341,1.9670) -- (0.3030,2.0030) -- (0.2720,2.0389) -- (0.2409,2.0749) -- (0.2099,2.1108) -- (0.1788,2.1468) -- (0.1478,2.1827) -- (0.1167,2.2186) -- (0.0857,2.2546) -- (0.0546,2.2905);
\draw[robotframe!50,dashed] (0.6756,1.5716) -- (0.6440,1.6070) -- (0.6110,1.6413) -- (0.5769,1.6743) -- (0.5416,1.7061) -- (0.5052,1.7366) -- (0.4678,1.7658) -- (0.4293,1.7937) -- (0.3899,1.8202) -- (0.3495,1.8452) -- (0.3083,1.8688) -- (0.2663,1.8909) -- (0.2235,1.9115) -- (0.1800,1.9306) -- (0.1358,1.9481) -- (0.0910,1.9640) -- (0.0458,1.9783) -- (-0.0000,1.9910);
\node[robotframe!70] at (-0.0000,1.9910) {$\times$};
\draw[robotframe!60,->] (2.1396,2.0192) -- (2.1867,2.0252) -- (2.2340,2.0296) -- (2.2815,2.0323) -- (2.3289,2.0333) -- (2.3764,2.0326) -- (2.4239,2.0303) -- (2.4712,2.0262) -- (2.5183,2.0204) -- (2.5653,2.0130) -- (2.6119,2.0039) -- (2.6581,1.9932) -- (2.7040,1.9808) -- (2.7494,1.9668) -- (2.7942,1.9511) -- (2.8385,1.9339) -- (2.8821,1.9152) -- (2.9251,1.8948) -- (2.9672,1.8730) -- (3.0086,1.8497) -- (3.0491,1.8249);
\draw[robotframe!60,->] (2.1396,2.0192) -- (2.1866,2.0261) -- (2.2336,2.0330) -- (2.2806,2.0399) -- (2.3276,2.0468) -- (2.3746,2.0537) -- (2.4216,2.0606) -- (2.4686,2.0675) -- (2.5156,2.0744) -- (2.5626,2.0813) -- (2.6096,2.0882) -- (2.6566,2.0951) -- (2.7036,2.1020) -- (2.7506,2.1089) -- (2.7976,2.1158) -- (2.8446,2.1227) -- (2.8916,2.1296) -- (2.9386,2.1365) -- (2.9856,2.1434) -- (3.0326,2.1503) -- (3.0795,2.1572);
\draw[robotframe!60,->] (2.1396,2.0192) -- (2.1865,2.0269) -- (2.2330,2.0363) -- (2.2792,2.0474) -- (2.3250,2.0601) -- (2.3703,2.0744) -- (2.4151,2.0903) -- (2.4592,2.1078) -- (2.5027,2.1269) -- (2.5455,2.1475) -- (2.5875,2.1696) -- (2.6288,2.1932) -- (2.6691,2.2182) -- (2.7086,2.2447) -- (2.7470,2.2725) -- (2.7845,2.3018) -- (2.8209,2.3323) -- (2.8562,2.3641) -- (2.8903,2.3971) -- (2.9232,2.4314) -- (2.9549,2.4667);
\draw[robotframe!60,->] (1.9454,2.2905) -- (1.9771,2.3259) -- (2.0100,2.3601) -- (2.0441,2.3932) -- (2.0794,2.4250) -- (2.1158,2.4555) -- (2.1533,2.4847) -- (2.1917,2.5126) -- (2.2312,2.5391) -- (2.2715,2.5641) -- (2.3127,2.5877) -- (2.3548,2.6098) -- (2.3976,2.6304) -- (2.4411,2.6495) -- (2.4852,2.6670) -- (2.5300,2.6829) -- (2.5753,2.6972) -- (2.6211,2.7099) -- (2.6672,2.7209) -- (2.7138,2.7303) -- (2.7607,2.7381);
\draw[robotframe!60,->] (1.9454,2.2905) -- (1.9765,2.3265) -- (2.0075,2.3624) -- (2.0386,2.3984) -- (2.0696,2.4343) -- (2.1007,2.4703) -- (2.1317,2.5062) -- (2.1628,2.5421) -- (2.1938,2.5781) -- (2.2249,2.6140) -- (2.2559,2.6500) -- (2.2870,2.6859) -- (2.3180,2.7219) -- (2.3491,2.7578) -- (2.3801,2.7938) -- (2.4112,2.8297) -- (2.4422,2.8656) -- (2.4733,2.9016) -- (2.5043,2.9375) -- (2.5354,2.9735) -- (2.5664,3.0094);
\draw[robotframe!60,->] (1.9454,2.2905) -- (1.9758,2.3270) -- (2.0049,2.3646) -- (2.0326,2.4031) -- (2.0590,2.4427) -- (2.0839,2.4831) -- (2.1073,2.5244) -- (2.1293,2.5665) -- (2.1498,2.6094) -- (2.1687,2.6529) -- (2.1860,2.6972) -- (2.2018,2.7420) -- (2.2160,2.7873) -- (2.2285,2.8331) -- (2.2394,2.8793) -- (2.2487,2.9259) -- (2.2562,2.9728) -- (2.2622,3.0199) -- (2.2664,3.0672) -- (2.2689,3.1147) -- (2.2698,3.1622);
\draw[robotframe!60,->] (1.6487,2.4433) -- (1.6496,2.4908) -- (1.6521,2.5382) -- (1.6563,2.5855) -- (1.6622,2.6326) -- (1.6698,2.6795) -- (1.6791,2.7261) -- (1.6900,2.7723) -- (1.7025,2.8182) -- (1.7167,2.8635) -- (1.7324,2.9083) -- (1.7498,2.9525) -- (1.7687,2.9961) -- (1.7892,3.0389) -- (1.8111,3.0810) -- (1.8346,3.1223) -- (1.8595,3.1628) -- (1.8859,3.2023) -- (1.9136,3.2409) -- (1.9427,3.2784) -- (1.9731,3.3149);
\draw[robotframe!60,->] (1.6487,2.4433) -- (1.6487,2.4908) -- (1.6487,2.5383) -- (1.6487,2.5858) -- (1.6487,2.6333) -- (1.6487,2.6808) -- (1.6487,2.7283) -- (1.6487,2.7758) -- (1.6487,2.8233) -- (1.6487,2.8708) -- (1.6487,2.9183) -- (1.6487,2.9658) -- (1.6487,3.0133) -- (1.6487,3.0608) -- (1.6487,3.1083) -- (1.6487,3.1558) -- (1.6487,3.2033) -- (1.6487,3.2508) -- (1.6487,3.2983) -- (1.6487,3.3458) -- (1.6487,3.3933);
\draw[robotframe!60,->] (1.6487,2.4433) -- (1.6479,2.4908) -- (1.6453,2.5382) -- (1.6411,2.5855) -- (1.6352,2.6326) -- (1.6276,2.6795) -- (1.6184,2.7261) -- (1.6075,2.7723) -- (1.5949,2.8182) -- (1.5808,2.8635) -- (1.5650,2.9083) -- (1.5476,2.9525) -- (1.5287,2.9961) -- (1.5083,3.0389) -- (1.4863,3.0810) -- (1.4628,3.1223) -- (1.4379,3.1628) -- (1.4116,3.2023) -- (1.3839,3.2409) -- (1.3548,3.2784) -- (1.3244,3.3149);
\draw[robotframe!60,->] (1.3244,2.5216) -- (1.3560,2.5570) -- (1.3890,2.5913) -- (1.4231,2.6243) -- (1.4584,2.6561) -- (1.4948,2.6866) -- (1.5322,2.7158) -- (1.5707,2.7437) -- (1.6101,2.7702) -- (1.6505,2.7952) -- (1.6917,2.8188) -- (1.7337,2.8409) -- (1.7765,2.8615) -- (1.8200,2.8806) -- (1.8642,2.8981) -- (1.9090,2.9140) -- (1.9542,2.9283) -- (2.0000,2.9410) -- (2.0462,2.9521) -- (2.0928,2.9615) -- (2.1396,2.9692);
\draw[robotframe!60,->] (1.3244,2.5216) -- (1.3554,2.5576) -- (1.3865,2.5935) -- (1.4175,2.6295) -- (1.4486,2.6654) -- (1.4796,2.7014) -- (1.5107,2.7373) -- (1.5417,2.7732) -- (1.5728,2.8092) -- (1.6038,2.8451) -- (1.6349,2.8811) -- (1.6659,2.9170) -- (1.6970,2.9530) -- (1.7280,2.9889) -- (1.7591,3.0249) -- (1.7901,3.0608) -- (1.8212,3.0968) -- (1.8522,3.1327) -- (1.8833,3.1686) -- (1.9143,3.2046) -- (1.9454,3.2405);
\draw[robotframe!60,->] (1.3244,2.5216) -- (1.3548,2.5581) -- (1.3839,2.5957) -- (1.4116,2.6342) -- (1.4379,2.6738) -- (1.4628,2.7142) -- (1.4863,2.7555) -- (1.5083,2.7976) -- (1.5287,2.8405) -- (1.5476,2.8840) -- (1.5650,2.9283) -- (1.5808,2.9731) -- (1.5949,3.0184) -- (1.6075,3.0642) -- (1.6184,3.1104) -- (1.6276,3.1570) -- (1.6352,3.2039) -- (1.6411,3.2510) -- (1.6453,3.2984) -- (1.6479,3.3458) -- (1.6487,3.3933);
\draw[robotframe!60,->] (1.0000,2.6000) -- (1.0008,2.6475) -- (1.0034,2.6949) -- (1.0076,2.7422) -- (1.0135,2.7894) -- (1.0211,2.8362) -- (1.0303,2.8828) -- (1.0412,2.9291) -- (1.0538,2.9749) -- (1.0679,3.0202) -- (1.0837,3.0650) -- (1.1011,3.1092) -- (1.1200,3.1528) -- (1.1405,3.1957) -- (1.1624,3.2378) -- (1.1859,3.2791) -- (1.2108,3.3195) -- (1.2371,3.3590) -- (1.2649,3.3976) -- (1.2940,3.4351) -- (1.3244,3.4716);
\draw[robotframe!60,->] (1.0000,2.6000) -- (1.0000,2.6475) -- (1.0000,2.6950) -- (1.0000,2.7425) -- (1.0000,2.7900) -- (1.0000,2.8375) -- (1.0000,2.8850) -- (1.0000,2.9325) -- (1.0000,2.9800) -- (1.0000,3.0275) -- (1.0000,3.0750) -- (1.0000,3.1225) -- (1.0000,3.1700) -- (1.0000,3.2175) -- (1.0000,3.2650) -- (1.0000,3.3125) -- (1.0000,3.3600) -- (1.0000,3.4075) -- (1.0000,3.4550) -- (1.0000,3.5025) -- (1.0000,3.5500);
\draw[robotframe!60,->] (1.0000,2.6000) -- (0.9992,2.6475) -- (0.9966,2.6949) -- (0.9924,2.7422) -- (0.9865,2.7894) -- (0.9789,2.8362) -- (0.9697,2.8828) -- (0.9588,2.9291) -- (0.9462,2.9749) -- (0.9321,3.0202) -- (0.9163,3.0650) -- (0.8989,3.1092) -- (0.8800,3.1528) -- (0.8595,3.1957) -- (0.8376,3.2378) -- (0.8141,3.2791) -- (0.7892,3.3195) -- (0.7629,3.3590) -- (0.7351,3.3976) -- (0.7060,3.4351) -- (0.6756,3.4716);
\draw[robotframe!60,->] (0.6756,2.5216) -- (0.6452,2.5581) -- (0.6161,2.5957) -- (0.5884,2.6342) -- (0.5621,2.6738) -- (0.5372,2.7142) -- (0.5137,2.7555) -- (0.4917,2.7976) -- (0.4713,2.8405) -- (0.4524,2.8840) -- (0.4350,2.9283) -- (0.4192,2.9731) -- (0.4051,3.0184) -- (0.3925,3.0642) -- (0.3816,3.1104) -- (0.3724,3.1570) -- (0.3648,3.2039) -- (0.3589,3.2510) -- (0.3547,3.2984) -- (0.3521,3.3458) -- (0.3513,3.3933);
\draw[robotframe!60,->] (0.6756,2.5216) -- (0.6446,2.5576) -- (0.6135,2.5935) -- (0.5825,2.6295) -- (0.5514,2.6654) -- (0.5204,2.7014) -- (0.4893,2.7373) -- (0.4583,2.7732) -- (0.4272,2.8092) -- (0.3962,2.8451) -- (0.3651,2.8811) -- (0.3341,2.9170) -- (0.3030,2.9530) -- (0.2720,2.9889) -- (0.2409,3.0249) -- (0.2099,3.0608) -- (0.1788,3.0968) -- (0.1478,3.1327) -- (0.1167,3.1686) -- (0.0857,3.2046) -- (0.0546,3.2405);
\draw[robotframe!50,dashed] (0.6756,2.5216) -- (0.6440,2.5570) -- (0.6110,2.5913) -- (0.5769,2.6243) -- (0.5416,2.6561) -- (0.5052,2.6866) -- (0.4678,2.7158) -- (0.4293,2.7437) -- (0.3899,2.7702) -- (0.3495,2.7952) -- (0.3083,2.8188) -- (0.2663,2.8409) -- (0.2235,2.8615) -- (0.1800,2.8806) -- (0.1358,2.8981) -- (0.0910,2.9140) -- (0.0458,2.9283) -- (-0.0000,2.9410);
\node[robotframe!70] at (-0.0000,2.9410) {$\times$};
\draw[robotframe!60,->] (0.3513,2.4433) -- (0.3521,2.4908) -- (0.3547,2.5382) -- (0.3589,2.5855) -- (0.3648,2.6326) -- (0.3724,2.6795) -- (0.3816,2.7261) -- (0.3925,2.7723) -- (0.4051,2.8182) -- (0.4192,2.8635) -- (0.4350,2.9083) -- (0.4524,2.9525) -- (0.4713,2.9961) -- (0.4917,3.0389) -- (0.5137,3.0810) -- (0.5372,3.1223) -- (0.5621,3.1628) -- (0.5884,3.2023) -- (0.6161,3.2409) -- (0.6452,3.2784) -- (0.6756,3.3149);
\draw[robotframe!60,->] (0.3513,2.4433) -- (0.3513,2.4908) -- (0.3513,2.5383) -- (0.3513,2.5858) -- (0.3513,2.6333) -- (0.3513,2.6808) -- (0.3513,2.7283) -- (0.3513,2.7758) -- (0.3513,2.8233) -- (0.3513,2.8708) -- (0.3513,2.9183) -- (0.3513,2.9658) -- (0.3513,3.0133) -- (0.3513,3.0608) -- (0.3513,3.1083) -- (0.3513,3.1558) -- (0.3513,3.2033) -- (0.3513,3.2508) -- (0.3513,3.2983) -- (0.3513,3.3458) -- (0.3513,3.3933);
\draw[robotframe!60,->] (0.3513,2.4433) -- (0.3504,2.4908) -- (0.3479,2.5382) -- (0.3437,2.5855) -- (0.3378,2.6326) -- (0.3302,2.6795) -- (0.3209,2.7261) -- (0.3100,2.7723) -- (0.2975,2.8182) -- (0.2833,2.8635) -- (0.2676,2.9083) -- (0.2502,2.9525) -- (0.2313,2.9961) -- (0.2108,3.0389) -- (0.1889,3.0810) -- (0.1654,3.1223) -- (0.1405,3.1628) -- (0.1141,3.2023) -- (0.0864,3.2409) -- (0.0573,3.2784) -- (0.0269,3.3149);
\draw[robotframe!50,dashed] (0.0546,2.2905) -- (0.0242,2.3270) -- (-0.0049,2.3646);
\node[robotframe!70] at (-0.0049,2.3646) {$\times$};
\draw[robotframe!50,dashed] (0.0546,2.2905) -- (0.0235,2.3265) -- (-0.0075,2.3624);
\node[robotframe!70] at (-0.0075,2.3624) {$\times$};
\draw[robotframe!50,dashed] (0.0546,2.2905) -- (0.0229,2.3259) -- (-0.0100,2.3601);
\node[robotframe!70] at (-0.0100,2.3601) {$\times$};
\fill (1.0000,0.7000) circle (2pt);
\draw[black,very thick,->] (1.0000,0.7000) -- (1.0000,1.1500);
\node[below right] at (1.0000,0.7000) {$q_S$};
\draw[green!50!black!100,<-] (7.8000,3.2000) -- (7.7992,3.2475) -- (7.7966,3.2949) -- (7.7924,3.3422) -- (7.7865,3.3894) -- (7.7789,3.4362) -- (7.7697,3.4828) -- (7.7588,3.5291) -- (7.7462,3.5749) -- (7.7321,3.6202) -- (7.7163,3.6650) -- (7.6989,3.7092) -- (7.6800,3.7528) -- (7.6595,3.7957) -- (7.6376,3.8378) -- (7.6141,3.8791) -- (7.5892,3.9195) -- (7.5629,3.9590) -- (7.5351,3.9976) -- (7.5060,4.0351) -- (7.4756,4.0716);
\draw[green!50!black!100,<-] (7.8000,3.2000) -- (7.8000,3.2475) -- (7.8000,3.2950) -- (7.8000,3.3425) -- (7.8000,3.3900) -- (7.8000,3.4375) -- (7.8000,3.4850) -- (7.8000,3.5325) -- (7.8000,3.5800) -- (7.8000,3.6275) -- (7.8000,3.6750) -- (7.8000,3.7225) -- (7.8000,3.7700) -- (7.8000,3.8175) -- (7.8000,3.8650) -- (7.8000,3.9125) -- (7.8000,3.9600) -- (7.8000,4.0075) -- (7.8000,4.0550) -- (7.8000,4.1025) -- (7.8000,4.1500);
\draw[green!50!black!100,<-] (7.8000,3.2000) -- (7.8008,3.2475) -- (7.8034,3.2949) -- (7.8076,3.3422) -- (7.8135,3.3894) -- (7.8211,3.4362) -- (7.8303,3.4828) -- (7.8412,3.5291) -- (7.8538,3.5749) -- (7.8679,3.6202) -- (7.8837,3.6650) -- (7.9011,3.7092) -- (7.9200,3.7528) -- (7.9405,3.7957) -- (7.9624,3.8378) -- (7.9859,3.8791) -- (8.0108,3.9195) -- (8.0371,3.9590) -- (8.0649,3.9976) -- (8.0940,4.0351) -- (8.1244,4.0716);
\draw[green!50!black!80,<-] (7.4756,4.0716) -- (7.4440,4.1070) -- (7.4110,4.1413) -- (7.3769,4.1743) -- (7.3416,4.2061) -- (7.3052,4.2366) -- (7.2678,4.2658) -- (7.2293,4.2937) -- (7.1899,4.3202) -- (7.1495,4.3452) -- (7.1083,4.3688) -- (7.0663,4.3909) -- (7.0235,4.4115) -- (6.9800,4.4306) -- (6.9358,4.4481) -- (6.8910,4.4640) -- (6.8458,4.4783) -- (6.8000,4.4910) -- (6.7538,4.5021) -- (6.7072,4.5115) -- (6.6604,4.5192);
\draw[green!50!black!80,<-] (7.4756,4.0716) -- (7.4446,4.1076) -- (7.4135,4.1435) -- (7.3825,4.1795) -- (7.3514,4.2154) -- (7.3204,4.2514) -- (7.2893,4.2873) -- (7.2583,4.3232) -- (7.2272,4.3592) -- (7.1962,4.3951) -- (7.1651,4.4311) -- (7.1341,4.4670) -- (7.1030,4.5030) -- (7.0720,4.5389) -- (7.0409,4.5749) -- (7.0099,4.6108) -- (6.9788,4.6468) -- (6.9478,4.6827) -- (6.9167,4.7186) -- (6.8857,4.7546) -- (6.8546,4.7905);
\draw[green!50!black!80,<-] (7.4756,4.0716) -- (7.4452,4.1081) -- (7.4161,4.1457) -- (7.3884,4.1842) -- (7.3621,4.2238) -- (7.3372,4.2642) -- (7.3137,4.3055) -- (7.2917,4.3476) -- (7.2713,4.3905) -- (7.2524,4.4340) -- (7.2350,4.4783) -- (7.2192,4.5231) -- (7.2051,4.5684) -- (7.1925,4.6142) -- (7.1816,4.6604) -- (7.1724,4.7070) -- (7.1648,4.7539) -- (7.1589,4.8010) -- (7.1547,4.8484) -- (7.1521,4.8958) -- (7.1513,4.9433);
\draw[green!50!black!50,dashed] (7.8000,4.1500) -- (7.7992,4.1975) -- (7.7966,4.2449) -- (7.7924,4.2922) -- (7.7865,4.3394) -- (7.7789,4.3862) -- (7.7697,4.4328) -- (7.7588,4.4791) -- (7.7462,4.5249) -- (7.7321,4.5702) -- (7.7163,4.6150) -- (7.6989,4.6592) -- (7.6800,4.7028) -- (7.6595,4.7457) -- (7.6376,4.7878) -- (7.6141,4.8291) -- (7.5892,4.8695) -- (7.5629,4.9090) -- (7.5351,4.9476) -- (7.5060,4.9851) -- (7.4756,5.0216);
\node[green!50!black!70] at (7.4756,5.0216) {$\times$};
\draw[green!50!black!50,dashed] (7.8000,4.1500) -- (7.8000,4.1975) -- (7.8000,4.2450) -- (7.8000,4.2925) -- (7.8000,4.3400) -- (7.8000,4.3875) -- (7.8000,4.4350) -- (7.8000,4.4825) -- (7.8000,4.5300) -- (7.8000,4.5775) -- (7.8000,4.6250) -- (7.8000,4.6725) -- (7.8000,4.7200) -- (7.8000,4.7675) -- (7.8000,4.8150) -- (7.8000,4.8625) -- (7.8000,4.9100) -- (7.8000,4.9575) -- (7.8000,5.0050);
\node[green!50!black!70] at (7.8000,5.0050) {$\times$};
\draw[green!50!black!50,dashed] (7.8000,4.1500) -- (7.8008,4.1975) -- (7.8034,4.2449) -- (7.8076,4.2922) -- (7.8135,4.3394) -- (7.8211,4.3862) -- (7.8303,4.4328) -- (7.8412,4.4791) -- (7.8538,4.5249) -- (7.8679,4.5702) -- (7.8837,4.6150) -- (7.9011,4.6592) -- (7.9200,4.7028) -- (7.9405,4.7457) -- (7.9624,4.7878) -- (7.9859,4.8291) -- (8.0108,4.8695) -- (8.0371,4.9090) -- (8.0649,4.9476) -- (8.0940,4.9851) -- (8.1244,5.0216);
\node[green!50!black!70] at (8.1244,5.0216) {$\times$};
\draw[green!50!black!80,<-] (8.1244,4.0716) -- (8.1548,4.1081) -- (8.1839,4.1457) -- (8.2116,4.1842) -- (8.2379,4.2238) -- (8.2628,4.2642) -- (8.2863,4.3055) -- (8.3083,4.3476) -- (8.3287,4.3905) -- (8.3476,4.4340) -- (8.3650,4.4783) -- (8.3808,4.5231) -- (8.3949,4.5684) -- (8.4075,4.6142) -- (8.4184,4.6604) -- (8.4276,4.7070) -- (8.4352,4.7539) -- (8.4411,4.8010) -- (8.4453,4.8484) -- (8.4479,4.8958) -- (8.4487,4.9433);
\draw[green!50!black!80,<-] (8.1244,4.0716) -- (8.1554,4.1076) -- (8.1865,4.1435) -- (8.2175,4.1795) -- (8.2486,4.2154) -- (8.2796,4.2514) -- (8.3107,4.2873) -- (8.3417,4.3232) -- (8.3728,4.3592) -- (8.4038,4.3951) -- (8.4349,4.4311) -- (8.4659,4.4670) -- (8.4970,4.5030) -- (8.5280,4.5389) -- (8.5591,4.5749) -- (8.5901,4.6108) -- (8.6212,4.6468) -- (8.6522,4.6827) -- (8.6833,4.7186) -- (8.7143,4.7546) -- (8.7454,4.7905);
\draw[green!50!black!80,<-] (8.1244,4.0716) -- (8.1560,4.1070) -- (8.1890,4.1413) -- (8.2231,4.1743) -- (8.2584,4.2061) -- (8.2948,4.2366) -- (8.3322,4.2658) -- (8.3707,4.2937) -- (8.4101,4.3202) -- (8.4505,4.3452) -- (8.4917,4.3688) -- (8.5337,4.3909) -- (8.5765,4.4115) -- (8.6200,4.4306) -- (8.6642,4.4481) -- (8.7090,4.4640) -- (8.7542,4.4783) -- (8.8000,4.4910) -- (8.8462,4.5021) -- (8.8928,4.5115) -- (8.9396,4.5192);
\draw[green!50!black!60,<-] (6.6604,4.5192) -- (6.6133,4.5252) -- (6.5660,4.5296) -- (6.5185,4.5323) -- (6.4711,4.5333) -- (6.4236,4.5326) -- (6.3761,4.5303) -- (6.3288,4.5262) -- (6.2817,4.5204) -- (6.2347,4.5130) -- (6.1881,4.5039) -- (6.1419,4.4932) -- (6.0960,4.4808) -- (6.0506,4.4668) -- (6.0058,4.4511) -- (5.9615,4.4339) -- (5.9179,4.4152) -- (5.8749,4.3948) -- (5.8328,4.3730) -- (5.7914,4.3497) -- (5.7509,4.3249);
\draw[green!50!black!60,<-] (6.6604,4.5192) -- (6.6134,4.5261) -- (6.5664,4.5330) -- (6.5194,4.5399) -- (6.4724,4.5468) -- (6.4254,4.5537) -- (6.3784,4.5606) -- (6.3314,4.5675) -- (6.2844,4.5744) -- (6.2374,4.5813) -- (6.1904,4.5882) -- (6.1434,4.5951) -- (6.0964,4.6020) -- (6.0494,4.6089) -- (6.0024,4.6158) -- (5.9554,4.6227) -- (5.9084,4.6296) -- (5.8614,4.6365) -- (5.8144,4.6434) -- (5.7674,4.6503) -- (5.7205,4.6572);
\draw[green!50!black!60,<-] (6.6604,4.5192) -- (6.6135,4.5269) -- (6.5670,4.5363) -- (6.5208,4.5474) -- (6.4750,4.5601) -- (6.4297,4.5744) -- (6.3849,4.5903) -- (6.3408,4.6078) -- (6.2973,4.6269) -- (6.2545,4.6475) -- (6.2125,4.6696) -- (6.1712,4.6932) -- (6.1309,4.7182) -- (6.0914,4.7447) -- (6.0530,4.7725) -- (6.0155,4.8018) -- (5.9791,4.8323) -- (5.9438,4.8641) -- (5.9097,4.8971) -- (5.8768,4.9314) -- (5.8451,4.9667);
\draw[green!50!black!50,dashed] (6.8546,4.7905) -- (6.8229,4.8259) -- (6.7900,4.8601) -- (6.7559,4.8932) -- (6.7206,4.9250) -- (6.6842,4.9555) -- (6.6467,4.9847) -- (6.6083,5.0126);
\node[green!50!black!70] at (6.6083,5.0126) {$\times$};
\draw[green!50!black!50,dashed] (6.8546,4.7905) -- (6.8235,4.8265) -- (6.7925,4.8624) -- (6.7614,4.8984) -- (6.7304,4.9343) -- (6.6993,4.9703) -- (6.6683,5.0062);
\node[green!50!black!70] at (6.6683,5.0062) {$\times$};
\draw[green!50!black!50,dashed] (6.8546,4.7905) -- (6.8242,4.8270) -- (6.7951,4.8646) -- (6.7674,4.9031) -- (6.7410,4.9427) -- (6.7161,4.9831) -- (6.6927,5.0244);
\node[green!50!black!70] at (6.6927,5.0244) {$\times$};
\draw[green!50!black!50,dashed] (7.1513,4.9433) -- (7.1504,4.9908) -- (7.1479,5.0382);
\node[green!50!black!70] at (7.1479,5.0382) {$\times$};
\draw[green!50!black!50,dashed] (7.1513,4.9433) -- (7.1513,4.9908) -- (7.1513,5.0383);
\node[green!50!black!70] at (7.1513,5.0383) {$\times$};
\draw[green!50!black!50,dashed] (7.1513,4.9433) -- (7.1521,4.9908) -- (7.1547,5.0382);
\node[green!50!black!70] at (7.1547,5.0382) {$\times$};
\draw[green!50!black!50,dashed] (8.4487,4.9433) -- (8.4479,4.9908) -- (8.4453,5.0382);
\node[green!50!black!70] at (8.4453,5.0382) {$\times$};
\draw[green!50!black!50,dashed] (8.4487,4.9433) -- (8.4487,4.9908) -- (8.4487,5.0383);
\node[green!50!black!70] at (8.4487,5.0383) {$\times$};
\draw[green!50!black!50,dashed] (8.4487,4.9433) -- (8.4496,4.9908) -- (8.4521,5.0382);
\node[green!50!black!70] at (8.4521,5.0382) {$\times$};
\draw[green!50!black!50,dashed] (8.7454,4.7905) -- (8.7758,4.8270) -- (8.8049,4.8646) -- (8.8326,4.9031) -- (8.8590,4.9427) -- (8.8839,4.9831) -- (8.9073,5.0244);
\node[green!50!black!70] at (8.9073,5.0244) {$\times$};
\draw[green!50!black!50,dashed] (8.7454,4.7905) -- (8.7765,4.8265) -- (8.8075,4.8624) -- (8.8386,4.8984) -- (8.8696,4.9343) -- (8.9007,4.9703) -- (8.9317,5.0062);
\node[green!50!black!70] at (8.9317,5.0062) {$\times$};
\draw[green!50!black!50,dashed] (8.7454,4.7905) -- (8.7771,4.8259) -- (8.8100,4.8601) -- (8.8441,4.8932) -- (8.8794,4.9250) -- (8.9158,4.9555) -- (8.9533,4.9847) -- (8.9917,5.0126);
\node[green!50!black!70] at (8.9917,5.0126) {$\times$};
\draw[green!50!black!60,<-] (8.9396,4.5192) -- (8.9865,4.5269) -- (9.0330,4.5363) -- (9.0792,4.5474) -- (9.1250,4.5601) -- (9.1703,4.5744) -- (9.2151,4.5903) -- (9.2592,4.6078) -- (9.3027,4.6269) -- (9.3455,4.6475) -- (9.3875,4.6696) -- (9.4288,4.6932) -- (9.4691,4.7182) -- (9.5086,4.7447) -- (9.5470,4.7725) -- (9.5845,4.8018) -- (9.6209,4.8323) -- (9.6562,4.8641) -- (9.6903,4.8971) -- (9.7232,4.9314) -- (9.7549,4.9667);
\draw[green!50!black!60,<-] (8.9396,4.5192) -- (8.9866,4.5261) -- (9.0336,4.5330) -- (9.0806,4.5399) -- (9.1276,4.5468) -- (9.1746,4.5537) -- (9.2216,4.5606) -- (9.2686,4.5675) -- (9.3156,4.5744) -- (9.3626,4.5813) -- (9.4096,4.5882) -- (9.4566,4.5951) -- (9.5036,4.6020) -- (9.5506,4.6089) -- (9.5976,4.6158) -- (9.6446,4.6227) -- (9.6916,4.6296) -- (9.7386,4.6365) -- (9.7856,4.6434) -- (9.8326,4.6503) -- (9.8795,4.6572);
\draw[green!50!black!60,<-] (8.9396,4.5192) -- (8.9867,4.5252) -- (9.0340,4.5296) -- (9.0815,4.5323) -- (9.1289,4.5333) -- (9.1764,4.5326) -- (9.2239,4.5303) -- (9.2712,4.5262) -- (9.3183,4.5204) -- (9.3653,4.5130) -- (9.4119,4.5039) -- (9.4581,4.4932) -- (9.5040,4.4808) -- (9.5494,4.4668) -- (9.5942,4.4511) -- (9.6385,4.4339) -- (9.6821,4.4152) -- (9.7251,4.3948) -- (9.7672,4.3730) -- (9.8086,4.3497) -- (9.8491,4.3249);
\fill (7.8000,3.2000) circle (2pt);
\draw[black,very thick,->] (7.8000,3.2000) -- (7.8000,2.7500);
\node[below right] at (7.8000,3.2000) {$q_G$};
\node[anchor=west] at (0,-.5) {\textcolor{robotframe}{Forward successors}\hspace{2cm}\textcolor{green!50!black}{Predecessors of the goal}};
\end{tikzpicture}

\caption{Motion-primitive trees shown in the $(x,y)$ plane; nodes also
carry heading. The start tree (red) contains forward successors. The
goal tree (green) contains predecessors, with arrows pointing toward
the goal. Both use the same forward-moving unicycle model. Dashed
extensions leave the domain and are rejected. No connection between
the two trees is claimed.}
\label{fig:l13-primitives}
\end{figure}

Two nodes at the same position but with different headings are different
configurations. Likewise, two spatially close endpoints may not be
connectable by the allowed motions. A valid connection must match the
required pose and respect the controls and collision constraints. If
obstacles are represented in workspace rather than already in configuration
space, check the robot's footprint along the motion, not only its reference
point.

\subsection{Growing backward from a goal pose}
For bidirectional planning, the goal tree should represent motions that
can actually arrive at $q_G$. Its new nodes are \emph{predecessors}:
configurations from which an allowed forward primitive reaches an existing
goal-tree node. Therefore, motion arrows point toward the goal-tree root.

For the constant-input model above, a predecessor can be computed by
integrating the same equations backward in time for $\Delta t$. This is
a construction operation: the robot executes the resulting edge forward
from the predecessor to its parent. It does not require the robot to drive
backward. Simply growing a forward-only robot away from the goal and then
reversing those motions would generally violate its admissible controls.

The final planning pipeline remains the same: generate valid connections,
find a route through them, and recover the corresponding robot motion.
The representation changes from free straight segments to motions that
respect the vehicle's kinematics.

\smallskip
\textit{Further reading.} The standard RRT construction is described in
\href{https://lavalle.pl/papers/Lav98c.pdf}{S.~M.~LaValle,
\emph{Rapidly-Exploring Random Trees: A New Tool for Path Planning} (1998)}.
The \href{https://ompl.kavrakilab.org/planners.html}{Open Motion Planning
Library planner overview} distinguishes basic planners from optimal variants
such as PRM* and RRT*.
